\subsection{Criteria for the integral to belong to E}

PROPOSITION 7. --- Let E be a Hausdorff locally convex space, and \(\mathbf{f} \in \mathcal{K}(\mathrm{X};\mathrm{E})\) . If \(\mathbf{f}(\mathrm{X})\) is contained in a complete convex subset A of E, then \(\int \mathbf{f} d\mu \in \mathbf{E}\) .

Let K be the support of f, which is compact by hypothesis. Since f is zero on X - K, \(\mathbf{f}(X)\) is equal to \(\mathbf{f}(K)\) or to \(\mathbf{f}(K) \cup \{0\}\) , therefore is compact since f is continuous and E is Hausdorff; the closed convex envelope C of \(\mathbf{f}(X)\) in E is then precompact (for the uniform structure induced by that of E) (TVS, II, §4, No.~1, Prop. 3). But since C is a closed subset of the complete space A, C is complete and therefore compact; a fortiori, C is compact for the weakened topology \(\sigma(\mathrm{E}, \mathrm{E}')\) ; but since that topology is induced by \(\sigma(\mathrm{E}'*, \mathrm{E}')\) , C is the closed convex envelope of \(\mathbf{f}(X)\) in \(E'\) for the topology \(\sigma(\mathrm{E}'*, \mathrm{E}')\) ; the proof is therefore concluded by the Corollary of Prop. 4 of No.~2.

COROLLARY 1. --- Let E be a Hausdorff locally convex space; for every function \(\mathbf{f} \in \mathcal{K}(\mathrm{X};\mathrm{E})\) , \(\int \mathbf{f} d\mu\) belongs to the completion \(\widehat{\mathbf{E}}\) of E.

Since the duals of E and \(\widehat{E}\) are identical, it suffices to apply Prop. 7 while regarding f as taking its values in \(\widehat{E}\) .

COROLLARY 2. --- If E is a quasi-complete Hausdorff locally convex space, then \(\int \mathbf{f} d\mu \in \mathrm{E}\) for every function \(\mathbf{f} \in \mathcal{K}(\mathrm{X};\mathrm{E})\) .

As noted at the beginning of the proof of Prop. 7, \(\mathbf{f}(\mathrm{X})\) is compact and its closed convex envelope C in E is precompact, hence bounded; but since the set C is closed and bounded, it is complete by hypothesis, and it suffices to apply Prop. 7.

We shall see, in Ch. VI, §1, No.~2, other criteria for \(\int \mathbf{f} d\mu\) to belong to E, which apply in particular to the functions in \(\widetilde{\mathcal{K}}(\mathrm{X};\mathrm{E})\) and not just those in \(\mathcal{K}(\mathrm{X};\mathrm{E})\) .

Corollary 2 of Proposition 7 may be applied in the following two cases: \(1^{\circ}\) E is a Banach space; \(2^{\circ}\) E is the dual of a barreled Hausdorff locally convex space G, and E is equipped with an S-topology, where S is a covering of G by bounded subsets (TVS, III, §4, No.~2, Cor. 4 of Th. 1). For example, Cor. 2 of Prop. 7 can be applied when E is the weak dual of a Banach space, or a space of measures \(\mathcal{M}(\mathrm{Y};\mathbf{C})\) equipped with the vague topology.

If \(\mathbf{X} = \mathbf{R}\) , \(\mu\) is Lebesgue measure on \(\mathbf{R}\) , and \(\mathrm{E}\) is a Banach space, then the integral \(\int \mathbf{f} d\mu\) of a function in \(\mathcal{K}(\mathrm{X};\mathrm{E})\) is none other than the integral

\[
\int_ {- \infty} ^ {+ \infty} \mathbf {f} (x) d x
\]

defined in FRV, II, §2, No.~1; this follows from formula (1), and from FRV, I, §1, No.~2, Cor. of Prop. 2.

